% \vspace{-0.05in}
\section{Related Work}
\label{sec:relate}

\vspace{-0.1in}
\subsection{Retrieval-Augmented Generation with LLMs}\vspace{-0.05in}
% Large Language Models (LLMs)\cite{zhao2023survey}, while exceptional at generating coherent and contextually relevant text, often encounter limitations when tasked with retrieving accurate, factual information from extensive external knowledge bases. Retrieval-Augmented Generation (RAG) systems address these challenges by retrieving pertinent information from external sources, augmenting the LLM's input, and grounding its responses in factual, domain-specific knowledge \cite{ram2023context, fan2024survey}. Traditional naïve RAG approaches \cite{hyde, gao2023retrieval, chan2024rq} embed queries into a vector space to retrieve the nearest matching context vectors, but this method is constrained by the limited context window of LLMs. To overcome these constraints, more advanced RAG systems have been developed that incorporate pre-retrieval, retrieval, and post-retrieval strategies beyond simple vector similarity \cite{gao2023retrieval}. Our proposed \model introduces a more efficient approach by utilizing entities as low-level keywords and extracting high-level keywords from relationships, effectively covering different layers of information. This approach maintains retrieval quality while reducing both indexing and retrieval overhead. Moreover, \model facilitates efficient incremental updates, allowing the graph to expand with minimal cost and complexity.
Retrieval-Augmented Generation (RAG) systems enhance LLM inputs by retrieving relevant information from external sources, grounding responses in factual, domain-specific knowledge \cite{ram2023context, fan2024survey}. Current RAG approaches \cite{hyde, gao2023retrieval, chan2024rq, yu2024rankrag} typically embed queries in a vector space to find the nearest context vectors. However, many of these methods rely on fragmented text chunks and only retrieve the top-k contexts, limiting their ability to capture comprehensive global information needed for effective responses.

Although recent studies \cite{edge2024local} have explored using graph structures for knowledge representation, two key limitations persist. First, these approaches often lack the capability for dynamic updates and expansions of the knowledge graph, making it difficult to incorporate new information effectively. In contrast, our proposed model, \model, addresses these challenges by enabling the RAG system to quickly adapt to new information, ensuring the model's timeliness and accuracy. Additionally, existing methods often rely on brute-force searches for each generated community, which are inefficient for large-scale queries. Our \model\ framework overcomes this limitation by facilitating rapid retrieval of relevant information from the graph through our proposed dual-level retrieval paradigm, significantly enhancing both retrieval efficiency and response speed.

\vspace{-0.1in}
\subsection{Large Language Model for Graphs}\vspace{-0.05in}
% Large Language Models (LLMs) have revolutionized natural language processing by excelling in tasks such as text understanding, reasoning, and generation. While LLMs are highly effective when working with sequential text data, there has been increasing interest in integrating LLMs with graph-structured data, which introduces richer, non-linear relationships between entities. This combination has led to significant advancements in various graph tasks, including node classification, edge prediction, and graph classification \cite{xia2024opengraph, tang2024graphgpt, tang2024higpt, guo2024graphedit, he2023explanations, yu2023empower}. For example, GraphGPT \cite{tang2024graphgpt} proposes a method that aligns LLMs with graph structural knowledge, leveraging LLMs’ ability to understand complex textual relationships and apply that understanding to graph structures. GraphEdit \cite{guo2024graphedit} utilizes LLMs to determine connectivity between nodes based on the textual information associated with each node, enhancing the model’s ability to infer relationships and enrich graph-based reasoning. Moreover, graph-based RAG systems represent an emerging area of research, where graphs are increasingly used as the primary indexing mechanism, allowing for more nuanced and structured retrieval. Advanced RAG systems such as KAPING \cite{baek2023knowledge} and G-Retriever \cite{he2024g} have integrated graph structures into their retrieval pipelines. GraphRAG \cite{edge2024local} employs graph communities to capture global information, using these communities to structure and retrieve large-scale knowledge for tasks requiring global reasoning.
Graphs are a powerful framework for representing complex relationships and find applications in numerous fields. As Large Language Models (LLMs) continue to evolve, researchers have increasingly focused on enhancing their capability to interpret graph-structured data. This body of work can be divided into three primary categories: i) \textbf{GNNs as Prefix} where Graph Neural Networks (GNNs) are utilized as the initial processing layer for graph data, generating structure-aware tokens that LLMs can use during inference. Notable examples include GraphGPT~\cite{tang2024graphgpt} and LLaGA~\cite{chenllaga}. ii) \textbf{LLMs as Prefix} involves LLMs processing graph data enriched with textual information to produce node embeddings or labels, ultimately refining the training process for GNNs, as demonstrated in systems like GALM~\cite{xie2023graph} and OFA~\cite{liuone}. iii) \textbf{LLMs-Graphs Integration} focuses on achieving a seamless interaction between LLMs and graph data, employing techniques such as fusion training and GNN alignment, and developing LLM-based agents capable of engaging with graph information directly~\cite{li2023grenade,brannon2023congrat}.